\section{工具与环境任务}
\label{sec:nlp-execution-environment}

程序任务的产物可以在固定副本上反复运行测试，环境中的一次修改却会留下来，影响后续所有操作。用户说“把我的预约改到通知中的时间和地点，不要发送通知，保留备注”，系统需要找到正确的记录，读取当前状态，执行允许的修改，再核验结果。工具返回、用户补充以及其他入口的修改都可能改变下一步该做什么，因此这类任务的基本单位不再是一次完整的文本生成，而是带观察的行动。本节把这样的系统称为\term{智能体}（agent）：语言模型在循环中依据目标和已有观察选择下一步，或调用工具，或回复用户；外部执行器负责真正执行，并把结果作为观察返回。

这一设定带来一串相互衔接的问题：模型先要获得调用工具的能力；工具很多时，要选出合适的工具并把参数绑定到可靠来源；多步任务要安排行动顺序，在观察出乎意料时调整，并让独立的子任务并行；轨迹变长后，要保存仍然有效的状态、取回相关历史，并把失败经验带到下一次尝试；提示驱动的能力不足时，要用交互轨迹训练模型；单个上下文难以兼顾执行与检查时，可以让多个模型实例分工；最后，执行失败要能恢复，完成与否要由环境证据判断。表\ref{tab:nlp-execution-agent-map}按这些挑战列出代表方法：前一行方法留下的缺口，往往就是后一行方法的出发点。

\begin{table*}[htbp]
\centering
\small
\begin{tabular}{p{18mm}p{30mm}p{48mm}p{44mm}}
\toprule
挑战 & 代表方法 & 核心机制 & 新增代价或局限 \\
\midrule
获得调用能力 & 函数调用格式、Toolformer、ToolLLM、Gorilla & 工具描述进入上下文，调用写成带特殊标记的结构化片段；自监督插入并按损失下降过滤；强模型合成指令与调用轨迹 & 合成数据有覆盖与正确性偏差；自监督只适合无副作用、可批量执行的工具 \\
选工具与填参数 & 描述嵌入检索、模式约束解码 & 按缺失信息检索候选工具，再查权限与副作用；参数逐项绑定来源，模式保证类型合法 & 合法的值仍可能指向错误对象 \\
规划 & Plan-and-Solve、ReWOO、ReAct、LLMCompiler & 先写计划再执行；每步依据最新观察决定；按依赖图并行派发 & 先规划难以应对意外；逐步交替调用多、上下文增长；并行依赖正确声明依赖 \\
记忆 & 结构化状态、Generative Agents、MemGPT & 按事件更新带来源与版本的状态；按新近度、重要性、相关性打分取回；分层上下文由模型自行换入换出 & 摘要可能丢失否定与未知项；检索命中相似但不适用的记录 \\
反思 & Reflexion & 把失败轨迹与评价写成文字教训，作为下一回合的上下文 & 不改参数；反思可能错误归因 \\
训练 & FireAct、AgentTuning；多轮强化学习（Search-R1、ArCHer） & 成功轨迹微调并混合通用数据；整条轨迹以任务结果为奖励，观察词元不计梯度 & 只能模仿示范；奖励稀疏、信用分配难、环境重置昂贵 \\
协作 & 角色分工（CAMEL、MetaGPT）、多智能体辩论 & 不同角色处理规划、执行与审查；多个实例互读答案后修订并聚合 & 通信与调用成本成倍增加；错误可在实例间传播 \\
恢复与完成 & 幂等键、独立回读、终态判定、pass\textasciicircum{}$k$ & 按失败类型分别恢复；完成以环境读取为准；重复运行衡量稳定性 & 依赖服务端契约；环境不能无损重置时重复试验不独立 \\
\bottomrule
\end{tabular}
\caption{工具与环境任务的挑战与方法谱系。每行列出一类挑战的代表方法、它们改变的机制以及随之带来的新代价；后续各小节按表中顺序展开。}
\label{tab:nlp-execution-agent-map}
\end{table*}

为了让各方法在同一任务上比较，本节使用一组教学业务状态。用户\code{U7}有个人预约\code{R17}，关联活动\code{ML-0418}，原时间为2025年4月17日14:00，地点字段为海棠楼101，备注为“带笔记”，版本为1。同一用户另有同名预约\code{R18}，关联其他活动，必须保持不变。共同通知给出的目标为4月18日周五14:00、海棠楼201；修改个人预约不修改活动通知，也不办理报名。模拟环境以\code{Haitang-201}表示房间，所有时间带北京时间偏移。配套程序是运行在本地内存中的确定性模拟，没有调用真实预约系统；以下轨迹、调用数与数值均为教学构造，不代表任何方法的实测效果。

\subsection{函数调用能力的获得}
\label{sec:nlp-execution-function-calling}

模型要使用工具，须先知道有哪些工具、调用应写成什么格式，以及在什么时候调用。最直接的做法是把工具描述写进提示，再用含调用的示范做有监督微调，但示范要人工编写，覆盖的工具有限。Toolformer让模型在普通文本上自己尝试插入调用，只保留确实有助于预测后文的那些；ToolLLM与Gorilla用强模型大规模合成指令与调用轨迹，把覆盖范围扩大到成千上万个真实接口。三条路线都要先约定调用在词元序列中的表示。

\subsubsection{工具描述与调用格式}
\label{sec:nlp-execution-call-format}

函数模式的写法（函数名、用途、参数的JSON模式）以及执行器校验参数、检查权限、注入身份的“调用—结果—回复”循环，已在第\ref{sec:nlp-dialogue-llm-function}小节给出。这里只补模型一侧：一次调用在词元序列中长什么样，模型怎样学会写出它。代码\ref{code:nlp-execution-update-schema}是教学目录中更新工具的描述，另附返回值说明；用途说明中的否定条件“不发送任何通知”决定了它能否用于本例。

\begin{codeblock}[label={code:nlp-execution-update-schema}]{更新工具的描述（教学示例）}
{
  "name": "update",
  "description": "修改单条个人预约的开始时间和房间；不发送任何通知，不改备注。",
  "parameters": {
    "type": "object",
    "properties": {
      "id":               {"type": "string"},
      "expected_version": {"type": "integer", "minimum": 1},
      "start":            {"type": "string", "format": "date-time"},
      "room":             {"type": "string"},
      "request_id":       {"type": "string"}
    },
    "required": ["id", "expected_version", "start", "room", "request_id"],
    "additionalProperties": false
  },
  "returns": "回执：status、new_version；版本不符时返回 CONFLICT"
}
\end{codeblock}

模型看到和写出的都是词元序列，因此工具描述、调用与结果必须按固定的\term{对话模板}（chat template）序列化。常见做法是：系统消息中列出全部工具描述；模型要调用工具时，在助手消息中输出一段由特殊标记包裹的结构化调用；执行器把结果写成一条工具消息追加到序列末尾，模型再继续生成。代码\ref{code:nlp-execution-call-template}给出一种教学模板，其中\code{<|assistant|>}、\code{<tool\_call>}等是加入词表的特殊词元。不同模型的模板写法不同，有的用标签包裹JSON，有的直接生成类似Python的函数调用；训练与推断必须使用同一套模板。

\begin{codeblock}[label={code:nlp-execution-call-template}]{含并行调用的对话序列化（教学模板）}
<|system|>可用工具：[{"name":"search",...},{"name":"read",...},
  {"name":"availability",...},{"name":"update",...}]<|end|>
<|user|>把我的读书会预约改到通知中的时间和地点，不要通知，保留备注。<|end|>
<|assistant|><tool_call>{"id":"c1","name":"search",
  "arguments":{"user":"U7","event":"ML-0418"}}</tool_call><|end|>
<|tool|>{"id":"c1","result":[{"id":"R17","title":"读书会"}]}<|end|>
<|assistant|><tool_call>{"id":"c2","name":"read","arguments":{"id":"R17"}}</tool_call>
<tool_call>{"id":"c3","name":"availability",
  "arguments":{"start":"2025-04-18T14:00:00+08:00","room":"Haitang-201"}}</tool_call><|end|>
<|tool|>{"id":"c3","result":{"available":true}}<|end|>
<|tool|>{"id":"c2","result":{"version":1,"room":"Haitang-101","note":"带笔记"}}<|end|>
\end{codeblock}

这个模板表达了两种多次调用。第一种是顺序调用：搜索返回编号以后，模型才能读取\code{R17}，于是调用与结果在助手消息和工具消息之间交替。第二种是\term{并行函数调用}（parallel function calling）：读取记录与查询时段互不依赖，模型在同一条助手消息中输出两段调用，各带一个调用编号；执行器可以同时执行，结果按完成先后返回，模型依靠编号而不是顺序把结果对应回调用。若没有编号，第二条工具消息先到时，模型就可能把时段查询结果当成记录内容。

运行时在解码过程中监视特殊词元：一旦出现\code{<tool\_call>}，就缓存后续词元直到\code{</tool\_call>}，再交给算法\ref{alg:nlp-dialogue-function-loop}那样的校验与执行流程，不合格的调用以错误消息返回给模型修正。也可以在调用片段内部启用约束解码，按JSON模式屏蔽非法词元（第\ref{sec:nlp-inference-grammar}小节），从源头保证片段可解析；参数值是否指向正确对象仍要另行检查（第\ref{sec:nlp-execution-tool-arguments}小节）。

有了模板，函数调用能力就可以用有监督微调获得：把多轮示范按模板序列化，在助手消息的词元上计算损失，包括思考文字、调用片段和最终回复；系统消息、用户消息和工具消息的词元则用式\eqref{eq:nlp-sft-masked-loss}的损失掩码排除。工具消息必须排除，因为它的内容由环境产生：若训练模型去预测工具结果，模型学到的恰是在没有真实返回时编造结果。示范从哪里来，是接下来两种方法要回答的问题。

\subsubsection{自监督工具调用学习：Toolformer}
\label{sec:nlp-execution-toolformer}

人工示范要回答两个难题：在一段文本的哪个位置调用工具才有帮助，以及调用写成什么参数。Toolformer的想法是让语言模型自己提出候选调用，执行后检查结果能否帮助预测后文；能帮助的调用保留下来，并入原文本，再用这些文本微调同一个模型\scite{nlp-execution-schick2023toolformer}。这样得到的训练数据只依赖少量调用示例和工具本身，不需要人工标注调用位置。

先约定调用的文本表示。一次调用$c=(a_c,i_c)$由工具名$a_c$和输入$i_c$组成，结果为$r$。用特殊标记把它线性化为
\[
 e(c)=\langle\mathrm{API}\rangle\,a_c(i_c)\,\langle/\mathrm{API}\rangle,\qquad
 e(c,r)=\langle\mathrm{API}\rangle\,a_c(i_c)\to r\,\langle/\mathrm{API}\rangle,
\]
前者只有调用，后者附上结果。原文使用的工具包括问答系统、计算器、维基百科搜索、机器翻译和日历，基座模型为GPT-J。

给定一段文本$\bm{x}=(x_1,\ldots,x_n)$和提示$P(\bm{x})$（提示中含该工具的若干调用示例，并要求模型在文本中插入调用），构造过程分四步。

第一步是采样调用位置。对每个位置$i$，计算模型在此处开始一次调用的概率
\[
 p_i=p_M\bigl(\langle\mathrm{API}\rangle\mid P(\bm{x}),x_{1:i-1}\bigr),
\]
只保留$p_i>\tau_s$的位置，若超过$k$个则取概率最高的$k$个。$\tau_s$是采样阈值。这一步用模型自己的判断筛掉大量不需要调用的位置。

第二步是采样调用。对每个保留位置$i$，以$[P(\bm{x}),x_1,\ldots,x_{i-1},\langle\mathrm{API}\rangle]$为前缀，最多采样$m$个调用$c_i^1,\ldots,c_i^m$，以$\langle/\mathrm{API}\rangle$为结束标记，没有生成结束标记的样本丢弃。

第三步是执行全部候选调用，得到结果$r_i$。工具必须能被批量自动执行，而且执行没有副作用。

第四步是按损失过滤。设$z$为放在文本之前的前缀，位置$i$之后的加权交叉熵为
\begin{equation}
 L_i(z)=-\sum_{j=i}^{n}w_{j-i}\log p_M(x_j\mid z,x_{1:j-1}),
 \label{eq:nlp-execution-toolformer-loss}
\end{equation}
其中权重$w_t$随距离$t$递减，原文取$\tilde w_t=\max(0,1-0.2t)$再归一化，只有紧随其后的五个词元计入。调用是否有用，比较三种前缀下的损失：
\begin{equation}
 L_i^+=L_i\bigl(e(c_i,r_i)\bigr),\qquad
 L_i^-=\min\Bigl(L_i(\varepsilon),\,L_i\bigl(e(c_i,\varepsilon)\bigr)\Bigr),\qquad
 L_i^--L_i^+\geq\tau_f ,
 \label{eq:nlp-execution-toolformer-filter}
\end{equation}
其中$\varepsilon$为空序列，$e(c_i,\varepsilon)$表示只写调用、不给结果。只有满足不等式的调用才保留，$\tau_f$为过滤阈值。$L_i^-$取两者中较小的一个，是为了把收益归于结果$r_i$：调用的输入本身可能已经透露了后文，例如输入中出现了要预测的词；若只与“不调用”比较，这种调用也会被误认为有用。与“只写调用”的损失也作比较，就要求结果在调用输入之外再提供信息。还有一处实现细节：调用与结果作为前缀放在文本之前，而不是插入位置$i$，因为此时模型尚未见过含调用的文本，插在中间会打断它熟悉的文本分布，使损失失去可比性。

过滤后，把保留的调用插回原位置得到$\bm{x}^*=(x_{1:i-1},e(c_i,r_i),x_{i:n})$，汇成增广语料，用普通的语言模型目标微调$M$。推断时，模型照常解码；一旦生成“$\to$”，运行时暂停解码，执行调用，把结果与$\langle/\mathrm{API}\rangle$写入序列，再继续生成。

下面用一个教学算例走一遍过滤准则，数值均为构造。文本为“读书会4月18日举行，距4月15日发布通知还有3天，请按时参加”，在“还有”之后的位置$i$采样到计算器调用$c$=\code{Calculator(18-15)}，结果$r=3$。五个权重为$\tilde w=(1,0.8,0.6,0.4,0.2)$，归一化后为$(1/3,4/15,1/5,2/15,1/15)$。设后续五个词元“3、天、，、请、按”在三种前缀下的负对数概率分别为
\[
\begin{aligned}
 \text{不调用：}&(2.40,0.30,0.20,1.10,0.90), & L_i(\varepsilon)&=1.127,\\
 \text{只写调用：}&(2.10,0.30,0.20,1.10,0.90), & L_i(e(c,\varepsilon))&=1.027,\\
 \text{调用加结果：}&(0.10,0.20,0.20,1.10,0.90), & L_i(e(c,r))&=0.333.
\end{aligned}
\]
以第一行为例，$2.40/3+0.30\times4/15+0.20/5+1.10\times2/15+0.90/15=0.800+0.080+0.040+0.147+0.060=1.127$。于是$L_i^-=1.027$，$L_i^--L_i^+=0.694$；取教学阈值$\tau_f=0.5$，这次调用保留。若在同一位置采样到问答调用，返回的是“海棠楼201”，与后文数字无关，设其带结果损失为$0.993$、只写调用损失为$1.060$，则$L_i^-=1.060$，差值仅$0.067$，被丢弃。

Toolformer与第\ref{sec:nlp-reasoning-rft}小节的拒绝采样微调同属“模型自己生成、按准则过滤、再微调”的自训练路线，区别在于过滤准则：拒绝采样按答案正误，Toolformer按语言模型损失的下降。增广语料除插入的调用外与预训练文本同源，原文借此让模型在学会调用的同时保留通用的语言建模能力。它的局限也来自这套构造：每个调用独立采样，模型学不到把一个工具的输出喂给另一个工具的链式调用；准则只看紧随其后的几个词元，只奖励能直接帮助预测后文的调用；工具必须在大规模文本上批量执行，因此只适用于没有副作用的查询类工具。本例的更新接口一旦执行就改动预约，无法用这种方式在预训练文本上试调用。

\subsubsection{工具使用指令数据的构造}
\label{sec:nlp-execution-tool-data}

Toolformer覆盖的是少量通用工具和单次调用。实际应用面对的是成千上万个接口，一条用户指令常要多次调用、组合多个接口。这需要大规模的“指令—调用轨迹”数据，人工编写不现实，于是改用强模型合成。

ToolLLM的数据构造分三步\scite{nlp-execution-qin2024toolllm}。第一步收集接口：从RapidAPI平台收集一万六千多个真实REST接口及其文档，按类别与集合组织。第二步生成指令：从接口集合中采样工具组合，分单工具、同类别多工具、同集合多工具三种情形，把采样接口的文档交给ChatGPT，要求它写出需要这些接口才能完成的指令，并列出相关接口；这些“指令—相关接口”对日后也用于训练接口检索器。第三步标注调用轨迹：让ChatGPT以函数调用方式真实调用接口，逐步完成指令。

第三步的困难在于，长轨迹中任何一步出错，ReAct式的单链生成就会整条失败，许多指令标注不出可用轨迹。ToolLLM为此提出\term{深度优先决策树搜索}（depth-first search-based decision tree, DFSDT）：把轨迹看成一棵树，每个节点是一个部分轨迹；在节点处由模型生成下一步调用作为子节点；若某条分支走到失败，或模型主动选择放弃，就回溯到父节点，生成一个与已试子节点不同的新子节点；找到第一条以有效答案结束的路径即停止。算法\ref{alg:nlp-execution-dfsdt}写出这一过程。

\begin{algorithm}[DFSDT调用轨迹标注]
\label{alg:nlp-execution-dfsdt}
\begin{algorithmic}[1]
\Require 指令$q$，接口集合$\mathcal{A}$及文档，标注模型$M$，调用预算$B$
\Ensure 一条以有效答案结束的轨迹，或标注失败
\Function{Expand}{$h$}\Comment{$h$为当前部分轨迹}
  \State $C\gets\emptyset$\Comment{本节点已试过的子节点}
  \While{$B>0$}
    \State 以$q,\mathcal{A},h$及已试子节点$C$为条件，令$M$生成与$C$不同的下一步$a$
    \If{$a$为给出答案}
      \State \Return $h\oplus a$
    \ElsIf{$a$为放弃}
      \State \Return 失败\Comment{回溯到父节点}
    \EndIf
    \State $o\gets$执行$a$；$B\gets B-1$；$C\gets C\cup\{a\}$
    \State $\tau\gets$\Call{Expand}{$h\oplus a\oplus o$}
    \If{$\tau$不是失败}
      \State \Return $\tau$
    \EndIf
  \EndWhile
  \State \Return 失败
\EndFunction
\State \Return \Call{Expand}{空轨迹}
\end{algorithmic}
\end{algorithm}

与ReAct式单链相比，DFSDT能从错误分支退回重试，提高了困难指令的标注成功率；与广度优先搜索相比，它只需要找到一条可行路径，不必展开同层所有分支，调用开销较小。生成新子节点时把已试子节点写进提示，是为了避免模型反复生成同一个失败调用。标注得到的轨迹用于微调LLaMA得到ToolLLaMA，评价时用模型评判的通过率与胜率衡量。

Gorilla面向的是机器学习模型库的接口调用，从TorchHub、TensorFlow Hub和HuggingFace收集接口文档，用自指令方法为每个接口生成指令\scite{nlp-execution-patil2024gorilla}。它的关键做法是\term{检索感知训练}（retriever-aware training）：训练时在指令后附上检索到的接口文档，写成“参考这份接口文档”的形式，让模型学会依据当前文档写调用，而不是只凭记忆。接口版本更新、参数改名后，测试时换入新文档即可，不必重新训练。评价时，Gorilla把生成的代码解析为抽象语法树，检查参考调用是否作为子树出现在其中；若生成的接口名不在接口库中，则记为幻觉。与运行代码相比，语法树匹配不需要执行环境，代价低，但只核对接口名与关键参数，不判断参数取值在语义上是否正确。

表\ref{tab:nlp-execution-tool-learning}对照三种获得调用能力的方式。合成数据的问题集中在两点：指令由文档反推，可能不像真实用户的需求；教师模型给出的轨迹可能本身有错，而按“得到答案”过滤也不保证答案正确。合成数据扩大了覆盖，但覆盖到的工具仍是文档描述清楚、能自动调用的那一部分。

\begin{table}[htbp]
\centering
\small
\begin{tabular}{p{17mm}p{24mm}p{28mm}p{26mm}}
\toprule
方法 & 监督来源 & 调用结构 & 过滤或评价 \\
\midrule
示范微调 & 人工或规则编写的多轮示范 & 任意，受示范覆盖限制 & 人工审核 \\
Toolformer & 模型在文本中自行插入调用 & 少量工具，单次、互不链接 & 后文损失下降$\geq\tau_f$ \\
ToolLLM & 强模型生成指令并搜索轨迹 & 大量真实接口，多步多工具 & DFSDT找到有效答案；模型评判 \\
Gorilla & 自指令生成指令，附检索文档 & 单次调用模型库接口 & 语法树子树匹配；幻觉率 \\
\bottomrule
\end{tabular}
\caption{获得函数调用能力的几种数据来源。各方法改变的是监督从哪里来、能覆盖怎样的调用结构，以及用什么准则筛选或评价数据。}
\label{tab:nlp-execution-tool-learning}
\end{table}

\subsection{工具检索与参数生成}
\label{sec:nlp-execution-tools}

模型具备调用能力之后，每一步仍要决定调用哪个工具、填什么参数。前者在工具很多时变成检索加过滤问题，后者要把每个值绑定到可靠来源；两者的错误都可能通过格式检查，需要按契约与对象身份另行核对。

\subsubsection{工具检索与动作选择}
\label{sec:nlp-execution-tool-selection}

教学目录有五个工具：只读的\code{search}支持按名称或按活动编号与用户查找，\code{read}读取单条记录，\code{availability}检查时段，\code{update}只改时间和房间且不发通知，\code{notify}向参与者发送消息。检索的查询应当是“当前还缺什么信息或变化”，而不是整段用户要求。尚无对象编号时应选搜索；已确认编号但不知道版本时应选读取；目标与前置条件都已确认，更新工具才成为候选。\code{notify}在语义上与“预约修改”高度相关，却违反本例的否定约束。

工具多时，常用描述嵌入做第一轮召回：以编码器$E$把查询$q$和工具描述$d_k$映射为向量，按余弦相似度$s(q,d_k)=\cos\bigl(E(q),E(d_k)\bigr)$取前$K$个候选写入上下文，再由模型在候选中选择。ToolLLM的接口检索器就用“指令—相关接口”对训练句向量编码器，以相关接口为正例、抽取的若干其他接口为负例做对比学习，目标与式\eqref{eq:nlp-matching-infonce}同类。召回之后还须过滤：调用者没有权限、缺少必需输入或副作用与约束冲突的工具，即使相似度最高也不能执行。模型生成目录外的工具名时，运行时应报告不可用，而不是拼凑一个接口。

名称相似的工具不能互换，需要比较的是契约。若真实系统的更新接口强制给所有参与者发消息，本例就没有满足约束的动作，正确做法是说明冲突或请求用户重新授权，而不是选一个“最接近”的工具。

\subsubsection{函数参数生成与模式约束}
\label{sec:nlp-execution-tool-arguments}

动作确定以后，参数要从有来源的字段绑定，而不是只把JSON填满。搜索结果提供\code{id}，刚读取的记录提供\code{expected\_version}，通知提供目标时间和房间，控制器生成任务内唯一的请求号。一个正常的更新请求为：

\begin{codeblock}[label={code:nlp-execution-tool-request}]{本地模拟工具的结构化参数}
{
  "id": "R17",
  "expected_version": 1,
  "start": "2025-04-18T14:00:00+08:00",
  "room": "Haitang-201",
  "request_id": "task7-goal1-v1"
}
\end{codeblock}

按代码\ref{code:nlp-execution-update-schema}的模式做约束解码，可以消除缺字段、多字段和类型错误，但对象绑定、时区与日历仍要逐项检查：解析带时区的时间，查询编号与版本，并把记录所属用户与活动编号同当前目标比对。本地接口核验中，\code{R404}得到不存在错误，字符串版本\code{"1"}得到模式错误，2月30日得到日期错误，过期版本得到冲突。更隐蔽的是两组合法候选：把日期填为通知发布日期4月15日，或把编号填为同名的\code{R18}。这些值都能通过模式与日期解析，若权限允许，接口会成功修改错误的目标。因此调用器在写前核对目标身份，结束判定再读取正确对象；参数文本与接口回执分别保存。

\subsection{计划与行动选择}
\label{sec:nlp-execution-planning}

多步任务要决定行动的顺序，以及观察出乎意料时怎样调整。预算至少包括模型调用次数、工具调用次数和总时间，失败调用与回查同样计入。三种控制结构在这些预算上各有取舍（图\ref{fig:nlp-execution-planning-patterns}）：先计划后执行只需很少的模型调用，计划也便于事先检查，但难以应对意外的观察；ReAct每一步都依据最新观察决定下一步，灵活，但模型调用多、上下文不断增长；依赖图让互不依赖的子任务并行，缩短等待时间。

\begin{figure*}[htbp]
  \centering
  % 第九卷 14.3.3：三种计划与行动组织方式（跨栏，约160 mm宽）。
% FigureModel：语言模型调用；FigureOutput：工具执行与调度；虚线：变量引用或重新规划。
\begin{tikzpicture}[
  x=1mm,y=1mm,
  figure picture,
  model/.style={figure node,draw=FigureModel,fill=FigureModelFill,
    minimum width=38mm,align=center,minimum height=11mm,inner xsep=2pt,line width=.8pt},
  tool/.style={figure node,draw=FigureOutput,fill=FigureOutputFill,
    minimum height=8mm,inner xsep=2pt},
  title/.style={font=\figurefont\small,text=FigureInk,align=center},
  ref/.style={figure flow,dash pattern=on 2pt off 1.5pt}
]
  % (a) ReAct
  \node[title] at (22,64) {（a）ReAct：逐步交替};
  \node[model] (ra) at (22,48) {语言模型\\Thought$_t$与Action$_t$};
  \node[tool,minimum width=30mm] (rt) at (22,22) {工具执行};
  \draw[figure flow] ([xshift=-7mm]ra.south) --
    node[figure edge label,left=1pt] {Action$_t$} ([xshift=-7mm]rt.north);
  \draw[figure flow] ([xshift=7mm]rt.north) --
    node[figure edge label,right=1pt] {Observation$_t$} ([xshift=7mm]ra.south);
  \node[figure note] at (22,7) {$T$次工具调用需$T{+}1$次模型调用\\每次重读提示与全部历史};

  % (b) ReWOO
  \node[title] at (74,64) {（b）ReWOO：先规划后执行};
  \node[model] (wp) at (74,48) {规划器\\一次写出含\#E$_k$的计划};
  \node[tool,minimum width=11mm] (w1) at (60,31) {\#E$_1$};
  \node[tool,minimum width=11mm] (w2) at (74,31) {\#E$_2$};
  \node[tool,minimum width=11mm] (w3) at (88,31) {\#E$_3$};
  \node[model] (ws) at (74,14) {求解器\\读计划与证据作答};
  \foreach \w in {w1,w2,w3} {
    \draw[figure flow] (wp.south -| \w) -- (\w.north);
    \draw[figure flow] (\w.south) -- (ws.north -| \w);
  }
  \draw[ref] (w1.east) -- (w2.west);
  \node[figure note] at (74,3) {工作器按变量引用代入结果\\共2次模型调用};

  % (c) LLMCompiler
  \node[title] at (126,64) {（c）LLMCompiler：依赖图并行};
  \node[model] (cp) at (126,48) {函数调用规划器\\输出任务依赖图};
  \node[tool,minimum width=34mm,align=center] (cf) at (126,34) {任务获取单元\\依赖满足即派发};
  \node[tool,minimum width=15mm] (c1) at (117,21) {执行器};
  \node[tool,minimum width=15mm] (c2) at (135,21) {执行器};
  \node[model] (cj) at (126,7) {汇合\\作答或重新规划};
  \draw[figure flow] (cp.south) -- (cf.north);
  \draw[figure flow] (cf.south -| c1) -- (c1.north);
  \draw[figure flow] (cf.south -| c2) -- (c2.north);
  \draw[figure flow] (c1.south) -- (cj.north -| c1);
  \draw[figure flow] (c2.south) -- (cj.north -| c2);
  \draw[ref] (cj.east) -- (148,7) -- node[figure edge label,right=1pt,text width=8mm] {重新\\规划} (148,48) -- (cp.east);
\end{tikzpicture}

  \caption{三种计划与行动的组织方式。（a）ReAct中模型每写出一个行动就等待真实观察，再决定下一步；（b）ReWOO由规划器一次写出含证据变量的计划，工作器依次执行并代入变量（虚线表示$\#E_2$引用$\#E_1$），求解器最后作答；（c）LLMCompiler的规划器输出任务依赖图，任务获取单元把依赖已满足的任务派给执行器并行执行，汇合步骤决定作答或重新规划。}
  \label{fig:nlp-execution-planning-patterns}
\end{figure*}

\subsubsection{计划与执行的分离}
\label{sec:nlp-execution-plan-execute}

零样本思维链（第\ref{sec:nlp-execution-cot}小节）常见的一类错误是漏掉步骤：模型直接写推理，写到一半忘了题目中的某个条件。Plan-and-Solve提示把一次推理拆成两段，先让模型理解问题并制订分步计划，再按计划逐步执行\scite{nlp-execution-wang2023planandsolve}；其增强版本还要求模型先抽取相关变量及其数值、注意中间结果的计算。它仍是一次生成，计划与执行写在同一段文字中，作用是让模型在执行前显式列出子目标，减少遗漏。

把这一思想用于工具调用，就得到ReWOO\scite{nlp-execution-xu2023rewoo}。它把系统拆成三个角色：\term{规划器}（planner）一次写出完整计划，每一步是一句说明加一个工具调用，调用结果存入证据变量$\#E_k$，后续步骤的输入可以引用前面的变量；\term{工作器}（worker）按顺序执行各步调用，把变量引用替换为实际结果；\term{求解器}（solver）读入任务、计划与全部证据，写出最终答案。规划时模型看不到任何观察，这正是方法名称“无观察推理”的含义。代码\ref{code:nlp-execution-rewoo-plan}是本例按此格式写出的计划。

\begin{codeblock}[label={code:nlp-execution-rewoo-plan}]{ReWOO格式的计划（教学示例）}
Plan: 按用户与活动中文名称查找预约。
#E1 = search[user="U7", name="机器学习读书会"]
Plan: 读取所找到记录的当前版本与字段。
#E2 = read[id=#E1.id]
Plan: 查询目标时段与房间是否可用。
#E3 = availability[start="2025-04-18T14:00:00+08:00", room="Haitang-201"]
Plan: 时段可用时按读取的版本更新时间和房间，不发通知。
#E4 = update[id=#E1.id, expected_version=#E2.version, ...]
\end{codeblock}

分离的收益首先体现在上下文开销上。设提示（任务说明、工具描述与示例）长$P$个词元，每一步的调用与观察平均增加$h$个词元，共$T$步。ReAct式交替要调用模型$T+1$次，第$t$次读入$P+(t-1)h$个词元；ReWOO只调用规划器和求解器各一次，求解器读入提示与全部证据。两者读入的词元总数分别约为
\begin{equation}
 C_{\mathrm{ReAct}}=(T+1)P+h\,\frac{T(T+1)}{2},\qquad
 C_{\mathrm{ReWOO}}=P_{\mathrm{plan}}+P_{\mathrm{solve}}+Th .
 \label{eq:nlp-execution-token-cost}
\end{equation}
前者随步数平方增长，后者线性增长。教学数值取$P=P_{\mathrm{plan}}=P_{\mathrm{solve}}=1500$、$h=200$、$T=5$，得$C_{\mathrm{ReAct}}=6\times1500+200\times15=12000$，$C_{\mathrm{ReWOO}}=1500+1500+1000=4000$。由于规划器不必读入观察，原文还把规划能力蒸馏到较小的模型上。

代价是计划无法根据观察调整。模拟目录存的是英文名称\code{ML reading group}，第一步按中文名称查询返回空列表，$\#E1$没有\code{id}字段，引用它的第二步和第四步都无法执行。纯粹的ReWOO只能把错误交给求解器，由它说明未完成。本例做一处扩展：工作器发现变量解析失败时，只把产生$\#E1$的子目标退回规划器，改按活动编号\code{ML-0418}查找，不依赖$\#E1$的$\#E3$及其结果保留。计划中含写操作时，写入前还须核对前面的证据确实满足条件，否则一次性计划会在版本已变或时段已占时照常提交。

\subsubsection{判断、行动与观察交替：ReAct}
\label{sec:nlp-execution-react}

先规划的方法在观察出乎意料时缺乏调整手段；只写思维链的方法没有外部信息，容易编造事实；只输出行动、不写推理的方法又难以分解目标、跟踪进度。ReAct把两者交织起来：在行动空间$\mathcal{A}$之外加入语言空间$\mathcal{L}$，模型在每一步可以写一段思考（thought），也可以发出一个行动\scite{yao2023react}。思考不改变环境，只改变上下文，用来分解目标、记录进度、解释观察和决定下一步。形式上，设$g$为目标，历史$h_t=(\mathit{th}_1,a_1,o_1,\ldots,\mathit{th}_{t-1},a_{t-1},o_{t-1})$，模型在每一步生成
\[
 (\mathit{th}_t,a_t)\sim\pi_{\bm{\theta}}(\,\cdot\mid g,h_t),\qquad a_t\in\mathcal{A},
\]
执行器执行$a_t$得到观察$o_t$。原文用少样本示范的轨迹作提示，示范逐行写成“Thought $t$”“Action $t$”“Observation $t$”三类行；在维基百科问答任务上，行动空间是\code{Search[实体]}、\code{Lookup[字符串]}和\code{Finish[答案]}。模型写完一行Action后，运行时以“Observation”为停止序列中断生成，执行行动，把真实返回写成Observation行，再让模型继续。Observation行永远由执行器写入，模型自己写出的“观察”不算数。代码\ref{code:nlp-execution-react-trace}是本例按这一格式记录的轨迹。

\begin{codeblock}[label={code:nlp-execution-react-trace}]{ReAct格式的执行轨迹（教学模拟）}
Question: 把U7在活动ML-0418的预约改到通知中的时间和地点，不要通知，保留备注。
Thought 1: 先要找到预约编号。按活动中文名称搜索。
Action 1: search[user="U7", name="机器学习读书会"]
Observation 1: []
Thought 2: 中文名称没有命中，改用确认过的活动编号。
Action 2: search[user="U7", event="ML-0418"]
Observation 2: [{"id": "R17", "event": "ML-0418"}]
Thought 3: 找到R17，读取当前版本与字段。
Action 3: read[id="R17"]
Observation 3: {"version": 1, "start": "2025-04-17T14:00:00+08:00",
               "room": "Haitang-101", "note": "带笔记"}
Thought 4: 修改前确认目标时段可用。
Action 4: availability[start="2025-04-18T14:00:00+08:00", room="Haitang-201"]
Observation 4: {"available": true}
Thought 5: 条件满足，按版本1更新时间和房间，不发通知。
Action 5: update[id="R17", expected_version=1, ...]
Observation 5: {"status": "committed", "new_version": 2}
Thought 6: 回执只说明请求已提交，还要读回字段。
Action 6: read[id="R17"]
Observation 6: {"version": 2, "start": "2025-04-18T14:00:00+08:00",
               "room": "Haitang-201", "note": "带笔记"}
Thought 7: 时间、地点已改，备注未变，任务完成。
Action 7: finish[R17已改到4月18日14:00海棠楼201，未发送通知，备注保留]
\end{codeblock}

这条轨迹共有六次工具调用。第1次的空结果使思考转向按编号查找，这正是ReAct相对一次性计划的优势：调整只发生在出问题的那一步。若把空结果误写成“已找到”，后续文字仍会很流畅，却已失去对象证据；若调用上限为4，执行器必须停在时段查询之后，报告尚未修改。配套程序中的动作由确定性规则选择，轨迹只展示格式与控制流，不能据此估计ReAct模型的成功率。

算法\ref{alg:nlp-execution-react}在这种交替上补充本章的预算与异常处理，这些工程条件不属于原论文设计。完成请求要先转成一次读取核验；写入结果未知时，下一步必须是查询原请求或读取对象，不允许另一个写动作越过它。

\begin{algorithm}[带预算与结果核验的行动循环]
\label{alg:nlp-execution-react}
\begin{algorithmic}[1]
\Require 目标$g$，工具目录，初始状态$S$，调用预算$B$，决策与时间上限
\Ensure 已核验完成，或附当前证据的未完成状态
\While{决策和时间预算未耗尽}
  \State 接收用户更正，更新目标版本，标记受影响的计划与事实
  \If{$B=0$}
    \State \Return 预算耗尽，附未核验项与结果未知的操作
  \EndIf
  \If{存在结果未知的写入}
    \State $a\gets$查询原请求状态或读取实际对象
  \Else
    \State 依据$g,S$写出思考并选择$a$；完成请求改为结束核验
  \EndIf
  \State 检查参数、身份、权限与前置条件；不符则有限次修正或退出
  \State $o\gets$实际调用$a$；$B\gets B-1$；保存参数、返回与时间
  \State 按事件更新$S$；写入超时记为未知，不记为成功或失败
  \If{最新核验观察满足$g$及保持条件，且无未知操作}
    \State \Return 已核验完成及其证据
  \EndIf
\EndWhile
\State \Return 达到上限，附已发生的变化与剩余条件
\end{algorithmic}
\end{algorithm}

ReAct与前后方法的关系可以这样概括：它是思维链加上外部行动，也是第\ref{sec:nlp-execution-swe-agent}小节SWE-agent在仓库修复中所用的基本循环。原文除提示外，还用模型自举得到的成功轨迹微调较小的模型，这条路线在第\ref{sec:nlp-execution-trajectory-sft}小节展开。它的代价已由式\eqref{eq:nlp-execution-token-cost}给出：模型调用次数与工具调用次数同阶，上下文随步数增长；模型还可能陷入重复同一行动的循环，或写出与随后行动不一致的思考，因此需要调用上限和重复检测。

\subsubsection{子目标依赖图与并行调用}
\label{sec:nlp-execution-subgoals}

ReAct每一步只发出一个行动，ReWOO的工作器按计划顺序执行，二者都把互不依赖的调用串行化了。工具调用的耗时常由网络与外部服务决定，串行等待会累积成主要延迟。LLMCompiler借用编译器的思路，把计划写成任务依赖图，再按依赖关系并行执行\scite{nlp-execution-kim2024llmcompiler}。它由三个部分组成：\term{函数调用规划器}（function calling planner）输出任务列表，每个任务是一次工具调用，输入可以用占位符\code{\$k}引用第$k$个任务的输出，占位符关系即依赖边；\term{任务获取单元}（task fetching unit）维护各任务状态，一旦某任务的全部前驱完成，就把占位符替换为实际输出并派发；\term{执行器}（executor）并行执行派来的调用。全部任务完成后，由模型汇合结果，决定直接作答还是重新规划；规划器还可以流式输出任务，前几个任务在计划写完前就开始执行。

设任务依赖图为$\mathcal{G}=(V,E)$，任务$v$的执行时间为$t_v$。顺序执行的总时间是$\sum_{v\in V}t_v$；按依赖并行时，总时间由图中耗时最长的路径决定，即\term{关键路径}（critical path）：
\begin{equation}
 T_{\mathrm{seq}}=\sum_{v\in V}t_v,\qquad
 T_{\mathrm{par}}=\max_{\text{路径}\,\rho\subseteq\mathcal{G}}\ \sum_{v\in\rho}t_v .
 \label{eq:nlp-execution-critical-path}
\end{equation}
并行缩短的是等待时间，调用次数不变。算法\ref{alg:nlp-execution-task-fetch}写出任务获取单元的调度。

\begin{algorithm}[任务获取单元的依赖调度]
\label{alg:nlp-execution-task-fetch}
\begin{algorithmic}[1]
\Require 任务依赖图$\mathcal{G}=(V,E)$，各任务的调用模板
\Ensure 各任务的输出，或失败任务及其后继
\State $R\gets\{v\in V:v\text{没有前驱}\}$；$D\gets\emptyset$\Comment{$R$为就绪集合，$D$为已完成集合}
\While{存在未完成任务}
  \For{$v\in R$}
    \State 把$v$输入中的占位符替换为前驱的实际输出，异步派发给执行器
  \EndFor
  \State $R\gets\emptyset$；等待任一任务$u$返回
  \If{$u$失败}
    \State 标记$u$在$\mathcal{G}$中的全部后继为不可执行，交汇合步骤处理
  \Else
    \State $D\gets D\cup\{u\}$
    \State $R\gets R\cup\{v:(u,v)\in E,\ v\text{的前驱全在}D\text{中}\}$
  \EndIf
\EndWhile
\end{algorithmic}
\end{algorithm}

为展示并行，另构造一个会前准备预约\code{R21}：用户明确允许4月18日13:00使用海棠楼202或203，优先202，读书会本身仍在201。任务图为

\begin{codeblock}[label={code:nlp-execution-compiler-plan}]{LLMCompiler格式的任务列表（教学示例）}
1. read(id="R21")
2. availability(room="Haitang-202", start="2025-04-18T13:00:00+08:00")
3. availability(room="Haitang-203", start="2025-04-18T13:00:00+08:00")
4. choose_room(prefer=["Haitang-202", "Haitang-203"], results=[$2, $3])
5. update(id="R21", expected_version=$1.version, room=$4, ...)
6. join()
\end{codeblock}

任务1、2、3没有前驱，可以同时派发；任务4依赖2与3，任务5依赖1与4。设读取耗时0.5秒，每次时段查询0.8秒，选择为本地计算不计时，更新0.6秒。顺序执行需$0.5+0.8+0.8+0.6=2.7$秒；并行时关键路径为“查询—选择—更新”，需$0.8+0.6=1.4$秒，两者都是四次服务调用。反过来，若202一查就可用，按优先规则的顺序策略可以不查203，只用三次调用；预先并行则多付一次调用，换取的是不必等待。并行是否划算，取决于延迟与调用成本哪个更紧。

依赖图还使更正的影响可以计算。用户把日期改为19日后，依赖日期的任务2、3及其后继4、5都要失效，按图的可达关系即可找出；任务1读取的记录与日期无关，结果保留。若只有任务2故障，也只需重做任务2及其后继。规划器的责任因此更重：漏写一条依赖边，例如让更新不依赖时段查询，执行器就可能在查询返回前提交写入。有副作用的任务不宜投机地提前并行，只读调用才适合积极并行。表\ref{tab:nlp-execution-planning-compare}汇总三种控制结构。

\begin{table}[htbp]
\centering
\small
\begin{tabular}{p{21mm}p{21mm}p{25mm}p{28mm}}
\toprule
方法 & 模型调用次数 & 对意外观察的响应 & 主要代价 \\
\midrule
ReWOO & 固定2次 & 计划内无法响应，交给求解器或重规划 & 变量解析失败使后续步骤作废 \\
ReAct & $T{+}1$次 & 每步依据最新观察调整 & 读入词元随步数平方增长，串行等待 \\
LLMCompiler & 规划与汇合各1次，可重规划 & 汇合时整体判断，可重规划 & 依赖声明错误导致竞争；可能多付调用 \\
\bottomrule
\end{tabular}
\caption{三种计划与行动结构的对照。$T$为工具调用步数；模型调用次数不含失败重试。}
\label{tab:nlp-execution-planning-compare}
\end{table}

\subsection{观察、状态与任务记忆}
\label{sec:nlp-execution-state-memory}

轨迹变长以后，执行器要有选择地保留信息：当前状态只保存仍影响任务的事实，上下文装不下的历史要压缩或按需取回，跨回合尝试还要带上失败经验。三个小节分别处理这三类信息。

\subsubsection{事件更新与结构化任务状态}
\label{sec:nlp-execution-state-update}

状态中的每个字段都应能回答“谁说的，针对哪个对象，在哪个版本成立”。把状态组织为$S=(g^{(v)},I,F,H,O)$：$g^{(v)}$是第$v$版用户目标，$I$为已确认的对象身份，$F$为带来源和版本的观察事实，$H$为尚未确认的假设，$O$为操作及其状态。每个事件$e_t$携带来源、对象编号、观察时间以及环境版本或请求号；更新$S_{t+1}=U(S_t,e_t)$按接口规则进行，不需要训练模型去猜写入是否发生。

正常轨迹中，第3次调用后$F$保存\code{R17}版本1的旧字段，$O$为“未提交”；第5次得到回执后，$O$记为“已提交，回执版本2”，但$F$中的字段仍待回读；第6次读取后$F$才有新快照。写入超时时，$O$记为“结果未知”，期望字段不复制进$F$。若其他入口在第5次与第6次之间改了备注，版本3的读取覆盖字段快照，版本2的回执留作审计；反过来，迟到的版本2响应不能覆盖已知的版本3。用户说“改到19日”更新的是$g$，不意味着环境已处于19日，两者分属不同字段。

\subsubsection{历史摘要与检索式记忆}
\label{sec:nlp-execution-history}

完整轨迹直接放进上下文能保留细节，却占用位置并混入过期信息。压缩历史有两条路：写成摘要，或存入外部记忆、按需检索。摘要至少要保存目标版本、对象编号、否定条件、已确认字段、未知操作及原始结果的定位。若摘要只写“预约改到18日”，就丢掉了“不发通知”和“保留备注”，执行器可能选择一个会通知全员的接口。较可靠的写法是“用户U7，活动ML-0418，目标18日14:00与201，不通知；R17版本1已读取；写入尚未提交”，并附各字段的来源事件号。遇到不确定的内容，摘要应保留“待核对”，不能在压缩时把“可能是R17”写成“已确认R17”。

检索式记忆要回答的问题是：记忆库很大时，当前该取回哪几条。Generative Agents把智能体经历的事件逐条写成自然语言记录，组成\term{记忆流}（memory stream），每条记录带创建时间与最近访问时间\scite{nlp-execution-park2023generative}。取回时，它对每条记忆$m$按三个分量打分：新近度按距上次访问的时间指数衰减，重要性由模型在1到10之间打分（日常琐事低、关键事件高），相关性是记忆与当前查询$q$的嵌入余弦相似度。三个分量各自在候选集合上做最小—最大归一化到$[0,1]$，再加权求和：
\begin{equation}
 \mathrm{score}(m,q)=\alpha_{\mathrm{rec}}\,\widehat{\gamma^{\Delta(m)}}
 +\alpha_{\mathrm{imp}}\,\widehat{\mathrm{imp}(m)}
 +\alpha_{\mathrm{rel}}\,\widehat{\cos\bigl(\bm{e}_m,\bm{e}_q\bigr)},
 \label{eq:nlp-execution-memory-score}
\end{equation}
其中$\Delta(m)$是距上次访问的时间，$\gamma<1$是衰减因子，帽子表示归一化，$\bm{e}_m,\bm{e}_q$为嵌入向量。原文以游戏内小时计时，取$\gamma=0.995$，三项权重均为1，按得分从高到低取能放进上下文的记忆。

用一个教学算例看三项如何相互作用。设当前查询为“更新R17的时间”，有三条候选记忆：$m_1$为“用户要求不要发送通知”，2小时前访问，重要性8，相关性0.62；$m_2$为“R18上次更新成功”，1小时前，重要性4，相关性0.71；$m_3$为“用户偏好下午活动”，30小时前，重要性3，相关性0.40。新近度为$0.995^2=0.990$、$0.995^1=0.995$、$0.995^{30}=0.860$，归一化后为$0.963$、$1$、$0$；重要性归一化后为$1$、$0.2$、$0$；相关性归一化后为$0.710$、$1$、$0$。总分分别为$2.673$、$2.200$、$0$。只按相关性，排第一的是另一个对象的$m_2$；加入重要性后，否定约束$m_1$排到第一。但$m_2$仍排第二，会被一起取回：打分只衡量文本是否值得读，不判断它是否适用于当前对象，取回后仍要按对象编号、时间与有效性过滤。原文还在重要性累计超过阈值时，让模型从近期记忆中归纳出更高层的\emph{反思}记录，也存入记忆流参与检索。

Generative Agents的检索规则是固定的：何时取、取多少由外部程序决定。MemGPT借用操作系统虚拟内存的思路，让模型自己管理记忆\scite{nlp-execution-packer2023memgpt}。它把上下文分为主上下文与外部上下文两层：主上下文相当于内存，含系统指令、一块可读写的工作上下文和一个先进先出的消息队列；外部上下文相当于磁盘，含保存全部历史消息的回忆存储和存放任意文本的归档存储。模型通过函数调用在两层之间搬运信息，例如把重要事实写入工作上下文、在归档存储中检索；队列接近上下文容量时，系统先向模型发出内存压力提示，让它保存要紧的信息，随后把较早的消息移出队列，并用递归摘要保留其要点。这种\term{虚拟上下文管理}（virtual context management）让取回决策随任务变化，代价是记忆操作本身也要消耗调用与上下文，模型漏存或误删时同样会丢信息。

两种机制都沿用第\ref{chap:nlp-classification-matching-organization}章的检索方法负责召回。评价记忆时，应同时检查关键条件的保留率、取回记录的身份正确率和最终执行结果；检索命中一条相似历史，只说明找到了文本，不说明当时的动作现在可以重用。

\subsubsection{反馈反思与回合记忆：Reflexion}
\label{sec:nlp-execution-reflexion}

一次尝试失败后，下一次尝试若只拿到原始轨迹，行动者要自己从长轨迹中找出关键错误。Reflexion让模型把失败写成一段简短的文字教训，作为下一回合的上下文\scite{nlp-execution-shinn2023reflexion}。它区分三个职责：\term{行动者}（actor）在环境中生成轨迹；评价器（evaluator）给出成败信号，可以是规则、单元测试、环境奖励或模型判断；\term{自我反思}（self-reflection）模型依据轨迹与评价写出反思。设第$j$次尝试的轨迹为$\tau_j$，评价为$r_j$，反思为
\[
 m_j=f_{\psi}(\tau_j,r_j),\qquad M_{j+1}=\mathrm{keep}(M_j\cup\{m_j\}),
\]
其中$\mathrm{keep}$只保留最近的少数几条。下一次尝试的行动者以原目标、新观察和$M_{j+1}$为条件。原文称之为“语言强化”：评价信号起到奖励的作用，但更新的是上下文中的记忆，而不是模型参数。评价通过或尝试次数用完时退出。

构造一条旧记忆：上一轮直接复用了历史编号\code{R404}，接口返回不存在。反思可以写成“上轮复用了历史编号，未核对当前活动；下一轮先按用户与活动编号重新查找，读取版本后再修改”，并链接到失败调用。下一轮只有真的先搜索、得到\code{R17}并完成核验，才能说反思起了作用。与另两种记忆条件相比：不保存记忆，容易再次复用旧编号；只保存原轨迹，要行动者自己定位错误；保存反思，提供了压缩后的修正规则，但反思也可能归因错误，例如把失败归于权限不足而建议换账号，因此外部错误码与反思器对原因的猜测应分别标注。

\subsection{智能体的训练}
\label{sec:nlp-execution-agent-training}

以上方法都靠提示驱动：基座模型能否写出合格的思考与调用，决定了智能体的上限。较小的模型按少样本ReAct提示工作时，常出现格式错误、重复行动和过早结束，长提示本身也占用上下文。训练可以直接提升多步工具使用。轨迹微调用成功轨迹做有监督学习，简单稳定，但只能模仿示范；多轮强化学习以任务结果为奖励在环境中探索，能超出示范，代价是奖励稀疏、轨迹长、环境交互昂贵。

\subsubsection{交互轨迹微调}
\label{sec:nlp-execution-trajectory-sft}

轨迹微调的数据流程是：用强模型在环境中按ReAct格式执行任务，记录完整轨迹，按任务结果过滤，再用保留的轨迹微调较小的模型。FireAct用GPT-4在问答任务上调用搜索接口生成轨迹，微调Llama等模型，并发现混合ReAct、思维链和Reflexion等不同格式的轨迹，能让模型按任务难度灵活选择策略\scite{nlp-execution-chen2023fireact}。AgentTuning从操作系统、数据库、知识图谱、网页购物等多个环境收集GPT-4轨迹，只保留奖励达标的轨迹组成智能体指令数据\scite{nlp-execution-zeng2024agenttuning}。

设一条轨迹为$\tau=(g,\mathit{th}_1,a_1,o_1,\ldots,\mathit{th}_T,a_T)$，序列化后为词元序列$(z_1,\ldots,z_{|\tau|})$。损失只计在模型应当生成的思考与行动词元上，目标、系统提示和观察词元用掩码$m_k$排除，与式\eqref{eq:nlp-sft-masked-loss}相同：
\begin{equation}
 L_{\mathrm{agent}}(\bm{\theta})=-\frac{1}{\sum_k m_k}\sum_{k=1}^{|\tau|}m_k\log p_{\bm{\theta}}(z_k\mid z_{<k}),\qquad
 m_k=\begin{cases}1,& z_k\in\mathit{th}_t\cup a_t,\\0,&\text{其他}.\end{cases}
 \label{eq:nlp-execution-trajectory-sft}
\end{equation}
观察词元仍留在序列中作为条件，只是不计损失；理由与工具消息相同，模型不应学习预测环境的返回。只用智能体轨迹微调，模型的通用能力会下降，AgentTuning因此采用\term{混合指令微调}（hybrid instruction tuning），以比例$\eta$混合智能体数据与通用对话数据：
\begin{equation}
 L(\bm{\theta})=\eta\,L_{\mathrm{agent}}(\bm{\theta})+(1-\eta)\,L_{\mathrm{general}}(\bm{\theta}),
 \label{eq:nlp-execution-hybrid-tuning}
\end{equation}
原文在7B模型上以0.1为步长搜索$\eta$，$\eta=0.2$在未见任务上效果最好。

按成功过滤再微调，与第\ref{sec:nlp-reasoning-rft}小节的拒绝采样微调同构，只是“答案正确”换成了“任务成功”；从强化学习的角度看，它是对成功示范的行为克隆。行为克隆的弱点也随之而来：训练只见过教师走过的成功状态，学生一旦偏离，例如第1次搜索返回了教师轨迹中没有出现的空结果，就进入训练数据没有覆盖的状态。过滤只保留成功轨迹，失败中的恢复动作恰好被丢掉了。这些问题促使研究者在环境中直接优化任务结果。

\subsubsection{面向多轮交互的强化学习}
\label{sec:nlp-execution-agent-rl}

把整条交互看成一个回合：状态是目标与截至当前的历史，行动是模型在一轮中生成的思考与调用，环境执行后返回观察，回合结束时由判定器给出奖励$R(\tau)$，例如任务成功记1、否则记0。目标是最大化期望奖励$J(\bm{\theta})=\mathbb{E}_{\tau\sim p_{\bm{\theta}}}[R(\tau)]$。这里用到策略梯度定理：期望奖励的梯度等于奖励加权的轨迹对数概率梯度的期望，$\nabla J=\mathbb{E}[R(\tau)\nabla\log p_{\bm{\theta}}(\tau)]$，减去不依赖行动的基线不改变期望，只降低方差；一般推导与PPO等算法可参考\BookChapterRef{chap:rl-learning}{卷七“未知模型的求解”章}。

多轮交互的特殊之处在于轨迹概率由两类因子组成：
\begin{equation}
 p_{\bm{\theta}}(\tau)=\prod_{t=1}^{T}\pi_{\bm{\theta}}(\mathit{th}_t,a_t\mid g,h_t)\;P_{\mathrm{env}}(o_t\mid g,h_t,a_t).
 \label{eq:nlp-execution-trajectory-prob}
\end{equation}
环境转移$P_{\mathrm{env}}$不含参数，取对数求梯度时整项消失，$\nabla\log p_{\bm{\theta}}(\tau)=\sum_t\nabla\log\pi_{\bm{\theta}}(\mathit{th}_t,a_t\mid g,h_t)$。展开到词元层面，只有模型生成的词元贡献梯度，观察词元不计入。这正是实现中“观察词元掩码”的来源：它不是经验技巧，而是策略梯度在含环境步的轨迹上的直接推论。若不加掩码，等于把环境返回当成模型的行动去强化，梯度方向就错了。

第\ref{sec:nlp-reasoning-grpo}小节的GRPO可以直接推广到多轮。对同一任务从相同初态采样$G$条轨迹$\tau_1,\ldots,\tau_G$，得到奖励$r_1,\ldots,r_G$，按式\eqref{eq:nlp-reasoning-group-advantage}计算组内标准化优势$\widehat A_i$，赋给第$i$条轨迹中所有模型生成的词元。与单轮目标式\eqref{eq:nlp-reasoning-grpo-objective}相比，唯一的改动是把对回答全部词元的平均换成只对模型生成词元的掩码平均：记$\ell_{i,k}(\bm{\theta})$为式\eqref{eq:nlp-reasoning-grpo-objective}中第$i$条轨迹第$k$个词元上的裁剪项与KL项之和，$m_{i,k}$为该词元的掩码，则
\begin{equation}
 J(\bm{\theta})=\mathbb{E}\Biggl[\frac{1}{G}\sum_{i=1}^{G}\frac{1}{\sum_k m_{i,k}}\sum_{k}m_{i,k}\,\ell_{i,k}(\bm{\theta})\Biggr].
 \label{eq:nlp-execution-multiturn-grpo}
\end{equation}
观察词元仍作为条件参与计算概率比，只是不进入求和。Search-R1是这一做法的实例：模型在推理中发出搜索请求，检索结果插回序列，训练时对检索结果词元做掩码，以最终答案是否与参考精确匹配为奖励，用PPO或GRPO更新\scite{nlp-execution-jin2025searchr1}。

与单轮推理相比，多轮强化学习新增了两个难点。其一是信用分配：一条十几步的轨迹只有末尾一个奖励，第1步改用编号搜索的正确决定与第5步的错误参数得到同一个优势。一种思路是引入轮次层面的价值估计，ArCHer在轮次层面用离策略的时序差分学习一个评价器，再用它给每一轮的词元级策略梯度提供优势\scite{nlp-execution-zhou2024archer}；另一种是为中间步骤设计过程奖励，原理见第\ref{sec:nlp-reasoning-prm}小节。其二是环境成本：组内比较要求从同一初态采样$G$条轨迹，就需要$G$个能重置到相同状态、互相隔离的沙箱；有副作用的真实服务无法这样使用，用户一方往往也要用模型模拟。奖励判定器还必须读取环境状态而不是轨迹中的完成声明，否则策略会学会只写“已完成”，判定方式见第\ref{sec:nlp-execution-recovery}小节。

\subsection{多智能体分工与协作}
\label{sec:nlp-execution-multi-agent}

单个智能体把规划、执行和检查放在同一个上下文中，有两个不便：上下文随任务增长，后期的步骤要在大量无关历史中找信息；同一个上下文检查自己的结果，容易沿用原来的错误判断，Huang等的实验表明，没有外部反馈时模型的自我纠正往往不能改进推理答案\scite{nlp-execution-huang2024selfcorrect}。多智能体方法把任务分给多个模型实例，每个实例有自己的角色提示与较短的上下文，彼此通过消息协作。主要有两类组织方式。

第一类是角色分工。CAMEL让两个实例分别扮演“用户”和“助手”，用户实例按任务逐步下达指令，助手实例执行并汇报，初始提示固定双方角色，防止对话中角色互换或偏离任务\scite{nlp-execution-li2023camel}。MetaGPT把软件开发的标准作业流程写进协作协议：产品经理、架构师、工程师和测试等角色依次产出需求文档、设计文档、代码和测试，角色之间交换结构化文档而非自由对话，各角色从共享消息池订阅与自己相关的消息\scite{nlp-execution-hong2024metagpt}。结构化消息减少了自然语言转述中的信息丢失，流程也规定了谁在什么时候交付什么。

第二类是辩论与聚合。Du等让$N$个实例独立回答同一问题，随后进行$R$轮辩论：每一轮中，每个实例读入其他实例上一轮的答案与理由，修订自己的答案；实例的答案通常在几轮后趋于一致，最后取多数答案\scite{nlp-execution-du2024debate}。记第$i$个实例在第$r$轮的回答为$y_i^{(r)}$，过程为
\[
 \begin{aligned}
 y_i^{(0)}&\sim\pi(\,\cdot\mid x),\qquad
 y_i^{(r)}\sim\pi\bigl(\,\cdot\mid x,\{y_j^{(r-1)}\}_{j=1}^{N}\bigr),\quad r=1,\ldots,R,\\
 \hat y&=\operatorname{majority}\bigl(y_1^{(R)},\ldots,y_N^{(R)}\bigr).
 \end{aligned}
\]
与式\eqref{eq:nlp-reasoning-majority-vote}的自一致性相比，辩论多了实例之间的信息交换：一个实例可以指出另一个实例的计算错误。代价是生成次数为$N(R+1)$，而且第$r$轮的每个上下文都要读入$N$份回答，上下文长度随$N$增长。实例之间互读也可能让错误传播：一个自信的错误答案说服其他实例，多数表决就会集中到错误上。比较时应在相同预算下进行，例如$N=3$、$R=2$共9次生成，应与9次采样的自一致性对照；Huang等在相同回答数下的比较中，辩论并未优于自一致性。

把分工用于本例，可以设规划者、执行者与审查者三个角色。规划者读用户要求，产出目标记录与子目标；执行者是唯一持有写权限的实例，按子目标调用工具；审查者只有只读权限，独立读取\code{R17}、\code{R18}与通知记录，对照目标逐项判定。消息用固定字段：发送者、接收者、消息类型（计划、执行报告、审查结论）、内容，以及所依据的调用编号。终止条件也写进协议：审查者给出通过结论，或协作轮数达到上限，或执行者报告无法满足的约束。单一写入者避免了两个实例重复提交，审查者不读执行者的完成声明，只读环境，相当于把第\ref{sec:nlp-execution-recovery}小节的独立判定做成了一个角色。

多智能体适合子任务需要不同提示、权限或上下文的情形，例如写代码与审代码、执行与审计；单个上下文就能容纳的短任务，增加的通信轮数与调用费用很难收回。判断收益时，应与“单智能体加自检”或“单智能体多次采样”在相同的调用与词元预算下比较。

\subsection{执行恢复与完成判断}
\label{sec:nlp-execution-recovery}

执行中断的原因各不相同，恢复路径也应分开。参数错误退回生成器修正字段；明确的临时故障按约定次数与退避间隔重试；版本冲突要重新读取、检查最新条件再重新规划；信息不足或存在多个候选对象时交给用户澄清，澄清方式见第\ref{chap:nlp-question-answering-dialogue}章。每条路径都设退出条件，不能把所有异常都接到“再试一次”。图\ref{fig:nlp-execution-environment-loop}把语言目标、控制状态与实际环境分开：观察只收录工具真正返回的内容，模型判断“应该已经改好”仍是待验证的假设。

\begin{figure*}[htbp]
  \centering
  % Adapted from academic-drawing/templates/tikz_template.tex.
% Main path uses south -> north ports; feedback is outside the main column.
\begin{tikzpicture}[
  every node/.style={font=\figurefont\small,align=center,text=BookInk},
  block/.style={book node,font=\figurefont\small,
    text width=43mm,minimum height=11mm},
  side/.style={block,text width=29mm},
  flow/.style={book arrow,rounded corners=1.2mm},
  recovery/.style={flow,dashed,draw=BookAmber}
]
  \node[block,draw=FigureInput,fill=FigureInput!10] (goal) at (0,0)
    {目标及保持条件\\对象、授权、目标版本};
  \node[side,draw=FigureInput,fill=FigureInput!10] (user) at (-61mm,0)
    {用户澄清\\或更正};
  \node[block,draw=FigureModel,fill=FigureModel!10] (plan) at (0,-24mm)
    {判断与计划\\依赖、预算、下一步};
  \node[block] (action) at (0,-48mm)
    {绑定参数并行动\\模式、身份、版本检查};
  \node[block,draw=FigureOutput,fill=FigureOutput!10] (environment)
    at (0,-72mm) {实际环境\\记录及副作用};
  \node[side] (external) at (61mm,-72mm)
    {其他入口\\修改共同状态};
  \node[block] (observation) at (0,-96mm)
    {真实工具观察\\字段、回执或异常};
  \node[block] (state) at (0,-120mm)
    {事件更新状态\\事实、假设、未知操作};
  \node[side,draw=BookAmber,fill=BookAmber!8] (recover) at (61mm,-120mm)
    {先回查旧请求\\恢复或重新读取};
  \node[block,draw=FigureOutput,fill=FigureOutput!10] (check) at (0,-144mm)
    {受控读取与结束核验\\目标及保持条件};
  \node[block,draw=FigureOutput,fill=FigureOutput!10] (done) at (0,-168mm)
    {输出完成证据\\或未完成与未知项};
  \draw[flow] (user.east) -- node[above] {更正} (goal.west);
  \draw[flow] (goal.south) -- (plan.north);
  \draw[flow] (plan.south) -- (action.north);
  \draw[flow] (action.south) -- (environment.north);
  \draw[flow] (external.west) -- (environment.east);
  \draw[flow] (environment.south) -- (observation.north);
  \draw[flow] (observation.south) -- (state.north);
  \draw[flow] (state.south) -- (check.north);
  \draw[flow] (check.south) -- (done.north);
  \draw[flow] (state.west) -- (-35mm,-120mm)
    -- (-35mm,-24mm) -- (plan.west);
  \node[text width=22mm,align=right,anchor=east] at (-39mm,-75mm)
    {更新信息\\重选行动};
  \draw[recovery] (state.east) --
    node[above,text width=17mm] {写入未知} (recover.west);
  \draw[recovery] (recover.north) -- (61mm,-96mm) -- (observation.east);
  \node[text width=29mm,align=center] at (61mm,-151mm)
    {回执描述已发生操作\\回读确认当前状态};
\end{tikzpicture}

  \caption{语言要求到环境结果的执行闭环。实线连接行动与真实观察；左侧反馈更新目标并重选行动，右侧虚线处理未知写入。用户更正影响目标，其他入口的修改影响环境，二者通过不同证据进入状态。}
  \label{fig:nlp-execution-environment-loop}
\end{figure*}
\FloatBarrier

写操作超时需要单独的状态。它可能尚未提交，也可能已提交而回执丢失，客户端看到的都是没有确定响应。控制器先记为“结果未知”，查询原请求的状态，再读取对象。模拟服务提供\code{request\_id}：相同请求号与相同参数的重放只返回原回执，实际写入一次；相同请求号换了参数则被拒绝。这称为\term{幂等键}（idempotency key），它依赖服务端持久保存请求与结果，客户端自己生成一个字符串并不能使重试变得安全。

三条教学轨迹核算了恢复的成本。写后超时：第5次调用已把版本1改为2但抛出超时，第6次查询请求得到已提交回执，第7次读取确认字段，共7次调用、1次写入。写前超时：第6次查询请求为不存在，本地服务同步执行、查询立即一致，在这一契约下可以排除请求仍在途；第7次以同一请求号和参数重试，第8次回读，共8次调用、1次写入。真实异步系统中“暂未查到”未必表示未执行，应继续有限回查或报告未知。第三条轨迹叠加了用户更正：第5次写入18日成功但回执超时，随后用户把个人预约改为19日15:00，并从另一入口把备注改成“带电脑”。状态先把目标更新为第2版，同时保留旧请求的未知状态；第6次查询确认旧请求已提交，第7次读到版本3，第8次检查19日时段，第9次以新请求号、版本3只改时间和房间，第10次读到版本4，确认19日15:00与“带电脑”。共10次调用、2次目标不同的写入；重用旧幂等键提交19日会被拒绝，盲目重发旧请求会恢复过时的18日目标，只把回复改成“已改到19日”则不改变任何记录。

完成判断应由独立的判定器读取环境，把“已完成”翻译成可检查的条件。设$g_j(s)$为第$j$项目标或保持条件在终态$s$上是否成立，终态判定为
\begin{equation}
 J(s)=\prod_{j=1}^{m}\mathbb{I}\bigl[g_j(s)\bigr].
 \label{eq:nlp-execution-completion}
\end{equation}
本例的条件是：\code{R17}的时间与地点为目标值，备注符合最新用户更正，\code{R18}不变，未发送通知。乘积要求全部成立，只检查时间就会放过地点仍旧的部分完成。条件取自任务规格与受控读取，不由被评价的执行器依据自己的回复生成；若过程另有要求，如修改前须获授权，则另查事件顺序，不能从终态反推过程合规。WebArena对信息查询比较答案，对导航与配置任务读取页面或存储状态后检查所需内容，不要求行动序列与参考轨迹一致\scite{nlp-intro-zhou2024webarena}；$\tau$-bench同样把对话结束时的数据库状态与标注的目标状态比较\scite{nlp-execution-yao2024taubench}，$\tau^2$-Bench进一步让用户一方也能改变环境\scite{nlp-intro-barres2025}。

表\ref{tab:nlp-execution-outcomes}摘要本地保存的运行记录。“预算为5”一行值得注意：环境已被正确修改，执行器却没有预算做第6次回读，只能报告“已提交，尚未核验”，独立判定随后确认目标达成。环境结果与执行器当时有权声称的结果是两个字段。

\begin{table}[htbp]
  \centering
  \begin{tabular}{lccc}
    \toprule
    轨迹 & 调用数 & 写入数 & 外部目标判定 \\
    \midrule
    正常执行 & 6 & 1 & 通过 \\
    只改时间、未改地点 & 1 & 1 & 不通过 \\
    读取服务不可用 & 3 & 0 & 不通过 \\
    预算为4，查询后停止 & 4 & 0 & 不通过 \\
    预算为5，未做回读 & 5 & 1 & 通过，未自验 \\
    成功修改同名另一对象 & 1 & 1 & 不通过 \\
    只宣布完成、未调用 & 0 & 0 & 不通过 \\
    \bottomrule
  \end{tabular}
  \caption{模拟执行的结果与成本。所有输入和轨迹均为教学构造；独立判定器的读取不计入执行器的调用预算。目标判定不隐藏环境故障或部分完成，停止原因另行记录。}
  \label{tab:nlp-execution-outcomes}
\end{table}

随机的智能体需要在重置后的环境中重复运行。式\eqref{eq:nlp-eval-passk}的pass@$k$衡量$k$次尝试中至少一次成功，适合可以独立重跑、由测试挑选的程序候选；环境任务通常只交付一次，更关心每次是否都成功。$\tau$-bench为此定义pass\textasciicircum{}$k$：对每个任务运行$n$次，其中成功$c$次，
\begin{equation}
 \text{pass\textasciicircum{}}k=\mathbb{E}_{\text{任务}}\left[\binom{c}{k}\Big/\binom{n}{k}\right],
 \label{eq:nlp-execution-pass-hat-k}
\end{equation}
即随机抽取的$k$次运行全部成功的概率。它随$k$增大而下降，下降越快说明行为越不稳定。例如某任务10次运行成功8次，pass\textasciicircum{}1为0.8，pass\textasciicircum{}2为$\binom{8}{2}/\binom{10}{2}=28/45\approx0.62$。两种指标都假定各次运行相互独立；若环境不能无损重置，第二次尝试已受第一次副作用影响，就不能按独立试验估计。成功比例之外，报告还应给出调用数、时延与重复副作用，并按工具选择、参数绑定、状态更新和终态判定分别统计失败。
